\section{The result and the plan}\label{sec:intro}

\subsection{Statement}
A frame of $n$ vectors in $\R^d$ is a matrix $X\in\R^{n\times d}$ whose $i$th
row is $x_i^T$. Distances are \emph{squared} Frobenius distances
$\fro{X-Y}^2=\sum_i\norm{x_i-y_i}^2$, with no normalisation. Throughout,
\[
 a=d/n .
\]
$X$ is \emph{Parseval} if $X^TX=I_d$ and \emph{equal-norm} if
$\norm{x_i}^2=a$ for all $i$; an equal-norm Parseval frame is an \emph{ENP frame}.
$X$ is an \emph{$\eps$-nearly ENP frame} if
\begin{equation}\label{eq:nearly}
 (1-\eps)I_d\preceq X^TX\preceq(1+\eps)I_d,
 \qquad (1-\eps)a\le\norm{x_i}^2\le(1+\eps)a\quad(1\le i\le n).
\end{equation}

\begin{theorem}[Linear Paulsen bound]\label{thm:main}
There is an absolute constant $C$ such that for all $1\le d\le n$, all
$0<\eps<1$ and every real $\eps$-nearly ENP frame $X\in\R^{n\times d}$ there is
an ENP frame $W\in\R^{n\times d}$ with $\fro{X-W}^2\le C\eps d$.
\end{theorem}

The bound is optimal up to the constant (\cref{rem:optimal}). Kwok, Lau, Lee and
Ramachandran~\cite{KLLR} proved the first bound independent of $n$,
$O(\eps d^{13/2})$; Hamilton and Moitra~\cite{HM} improved it to
$O(\eps d^2)$; Lau and Ramachandran announced the linear bound~\cite{LR}, via a
different route. \Cref{app:formal} describes a Lean~4 formalisation of
\cref{thm:main,thm:projection}.

For a Parseval $U$ write $P=UU^T$ for the orthogonal projection onto its column
space; then $P_{ii}=\norm{u_i}^2$. The problem has an equivalent form in which
the error is additive.

\begin{theorem}[Projection form]\label{thm:projection}
There is an absolute constant $C$ such that for every $n\ge1$ and every
rank-$d$ orthogonal projection $P$ on $\R^n$, with
$\beta=\max_i|P_{ii}-d/n|$, there is a rank-$d$ orthogonal projection $Q$ with
$Q_{ii}=d/n$ for all $i$ and $\fro{P-Q}^2\le Cn\beta$.
\end{theorem}

This reformulation and the complement symmetry are due to Cahill and
Casazza~\cite{CC}: both $\beta$ and $\fro{P-Q}$ are unchanged when $(P,Q)$ is
replaced by $(I-P,I-Q)$, so in \cref{thm:projection} one may assume $d\le n/2$.
In \cref{sec:assembly}, \cref{thm:projection} is derived from an equal-norm
formulation (\cref{prop:equalrow}) and \cref{thm:main} from
\cref{thm:projection}.

\begin{remark}[Optimality]\label{rem:optimal}
Let $d=2d_1$, $n=2n_1$, and $1\le k\le n_1/4$ with $n_1-k\ge d_1$; put
$\eps=4k/n$. Let $X=X_1\oplus X_2$, where $X_1$ is an ENP frame of $n_1-k$
vectors in $\R^{d_1}$ and $X_2$ one of $n_1+k$ vectors in $\R^{d_1}$
(\cref{cor:exist}). Then $X$ is Parseval, and its squared row norms
$a(1\mp2k/n)^{-1}$ make it an $\eps$-nearly ENP frame. Let $W$ be ENP, and let $E$
and $A$ be the principal blocks of $XX^T$ and $WW^T$ on the first $n_1-k$ rows.
Then $E$ is a projection of trace $d_1$ and $A$ a positive contraction of trace
$(n_1-k)a$, so $\Delta=E-A$ has $\tr\Delta=ka=\eps d/4$. Since
$\tr(E\Delta)=\tr\Delta+\tr((I-E)A)\ge\tr\Delta$,
\[
 \tr(A-A^2)=-\tr\Delta+2\tr(E\Delta)-\tr\Delta^2\ge\tr\Delta-\fro\Delta^2 .
\]
As $WW^T$ is a projection, its off-diagonal block has squared norm $\tr(A-A^2)$,
while that of $XX^T$ vanishes; hence
$\fro{XX^T-WW^T}^2\ge\fro\Delta^2+2(\tr\Delta-\fro\Delta^2)\ge2\tr\Delta-\fro{XX^T-WW^T}^2$,
i.e.\ $\fro{XX^T-WW^T}^2\ge\tr\Delta$, and by \cref{lem:align}
$\fro{X-W}^2\ge\frac12\tr\Delta=\eps d/8$.
\end{remark}

\paragraph{Conventions.} $c,C$ denote positive absolute constants whose value
may change between occurrences; constants with subscripts or names are fixed
once chosen. $\opn{\cdot}$ is the Euclidean operator norm, $\one$ the all-ones
vector, $\Diag(x)$ the diagonal matrix with diagonal $x$, $\diag(M)$ the diagonal
of $M$ as a vector, and $\Pi_{\mathcal S}$ the orthogonal projection onto a
subspace $\mathcal S$. $A\preceq B$ is the Loewner order. All Gaussian vectors are
real and centred.

\subsection{The plan}\label{subsec:plan}

\paragraph{The engine: scaling.}
Let $U$ be Parseval with projection $P$ and diagonal $p=\diag P$. For
$s\in\R^n$ the matrix $\Diag(e^s)U$ has a column space with projection $P(s)$,
and $s\mapsto\diag P(s)$ is the gradient of the convex function
$\frac12\log\det(U^T\Diag(e^{2s})U)$, whose Hessian at $s=0$ is $2L_P$, where
\[
 L_P=\Diag(p)-P\circ P
\]
is the graph Laplacian with edge weights $P_{ij}^2$ (\cref{lem:energy}). To
first order, changing the diagonal by $\delta\perp\one$ therefore requires
$s\approx\frac12L_P^{+}\delta$, at a cost of about
$s^TL_Ps\approx\frac14\delta^TL_P^{+}\delta$. If $L_P$ has spectral gap
$\lambda$ and $|\delta_i|\le\eps a$, the cost is at most $n\eps^2a^2/\lambda$.
For a ``random-looking'' $P$ one has $\lambda\approx a$ and the cost is
$\eps^2 d$, but $L_P$ can be arbitrarily badly connected: for an exact direct
sum the components cannot exchange diagonal mass at all.

\paragraph{Perturb, then scale.}
Perturb $U$ by a random horizontal tangent of size $t=\sqrt{K_0\eps}$, with
$K_0$ a large constant, at cost $O(t^2d)=O(\eps d)$. Generic entries of the new
projection then have size $t\sqrt{a/n}$, so most edges acquire weight
$\gtrsim at^2/n$, and the random walk with rates $P_{ij}^2$ reaches any half of
the vertices in time $O(1/(at^2))$: the \emph{barrier constant}
(\cref{def:barrier,lem:core}) is $O(1/(at^2))$. The static balancing theorem
(\cref{thm:balancing}) then removes exactly every diagonal error $\beta$ with
$\beta\cdot O(1/(at^2))\le\frac12$, at cost $\frac{15}4n\beta=O(t^2d)$. The
$\ell_\infty$ (rather than spectral) formulation is essential: the scaling must
stay in a box $\norm{s}_\infty\le1$, and an $\ell_2$ argument loses a factor
$\sqrt n$.

\paragraph{The difficulty.}
The perturbed frame must therefore have diagonal error at most $\theta at^2$ in
$\ell_\infty$, for a small fixed $\theta$. The initial error $\eps a=at^2/K_0$
is acceptable. But a random tangent of size $t$ moves each diagonal entry
\emph{linearly} by about $t\sqrt{a/n}$, which exceeds $at^2$ as soon as
$\eps\lesssim1/d$. Two different constructions deal with this.

\paragraph{Seed 1: many rows ($n\ge Bd^2$, \cref{sec:manyrow}).}
Here the input is an equal-norm frame $X$ with spectral error
$\eta=\opn{X^TX-I}$, and $t^2$ is a large multiple of $\eta$. Perturb each row
\emph{orthogonally to itself} and renormalise. Row norms stay exact, but the
frame operator moves. Its first-order change is removed by conditioning the
Gaussian on a subspace of codimension at most $d^2$, and its second-order
fluctuation is an average over $n$ rows, of size $t^2d/\sqrt n$. The spectral
error of the perturbed frame is therefore $\eta+O(t^2d/\sqrt n+t^2d^2/n+t^3)$,
which is at most $\theta t^2$ once $t^2/\eta$ and $B$ are large and $t$ is
small; whitening turns a spectral error $\delta$ into a diagonal error at most
$2a\delta$. The cubic term is $O(t^3)$, not the termwise $O(t^3d)$, because the
first-order constraint allows the normalisation weights to be centred
(\cref{lem:remainder}); this is why the seed works for every $\eta\le c_*$,
with no relation required between $\eta$ and $d$.

\paragraph{Seed 2: moderate rows ($d\ge A\log 2n$, \cref{sec:moderate}).}
Perturb $U$ in the horizontal tangent space of the Grassmannian, so that the
frame stays exactly Parseval, with a Gaussian whose covariance is filtered:
$C_\rho=\rho(I-\Omega)(\rho I+\Omega)^{-1}$, where $\Omega$ is the normalised
Gram operator of the linear diagonal map $\mathcal A$ and $\rho=D^2t^2$ for a
large constant $D$. The filter makes the linear diagonal change
$O(t\sqrt{\rho a\log n/n})=O(Dat^2\sqrt{\log n/d})$, negligible when
$d\gg\log n$, while leaving the weak (connecting) directions untouched. The
price is a second-order bias $t^2b_*$, which is not small in $\ell_\infty$. A
commutator identity shows that $b_*=\mathcal Am_*$ for an explicit
deterministic tangent direction $m_*$ with $\fro{m_*}^2\le2a/\rho$
(\cref{lem:commutator,lem:mean}). We add the drift $\frac t2m_*$ to the noise:
its first-order effect cancels the bias exactly, and its other effects are
$O(at^2/D)$. All remaining errors are at most $\theta at^2$ by concentration,
which is where $d\ge A\log 2n$ is used. Only a bounded, deterministic set of
rows needs a spectral argument, and second moments suffice there.

\paragraph{Assembly (\cref{sec:assembly}).}
After complementation and normalisation it suffices to treat equal-norm frames
with $2d\le n$ and spectral error $\eta$. If $d<d_0$ (a large absolute
constant), the Hamilton--Moitra bound $\eta d(d-1)\le d_0\eta d$ suffices
(\cref{sec:bounded}). If $d\ge d_0$ and $n\ge Bd^2$, Seed~1 applies. If
$d\ge d_0$ and $n<Bd^2$, then $A\log 2n\le A\log(2Bd^2)\le d$ and Seed~2
applies.

\begin{center}
\small
\begin{tabular}{@{}lll@{}}
\toprule
regime & construction & where\\
\midrule
$d<d_0$ & radial isotropic position, ordered coordinates (Hamilton--Moitra) & \cref{sec:bounded}\\
$d\ge d_0,\ n\ge Bd^2$ & row-tangent Gaussian, row normalisation, whitening, balancing & \cref{sec:manyrow}\\
$d\ge d_0,\ 2d\le n<Bd^2$ & filtered horizontal Gaussian with drift, balancing & \cref{sec:moderate}\\
$d>n/2$ & complement & \cref{sec:assembly}\\
\bottomrule
\end{tabular}
\end{center}
